\section{The semismall Springer decomposition}
\label{chapter_bbdg}

We recall the form of the decomposition theorem needed for the Springer
resolution. Let $f:X\to Y$ be a proper morphism of irreducible complex
algebraic varieties, and choose a stratification
$Y=\bigsqcup_\alpha Y_\alpha$ over which $f$ is topologically locally
trivial. The morphism $f$ is \vocab{semismall} if
\[
  \dim(X\times_YX)\leq\dim X.
\]
Equivalently, for every $y\in Y_\alpha$,
\[
  2\dim f^{-1}(y)\leq\operatorname{codim}_Y(Y_\alpha).
\]
A stratum is \vocab{relevant} when equality holds.

\begin{theorem}[thm:bbdg_semismall]{\cite{Beilinson-Bernstein-Deligne1982}}
  Let $f:X\to Y$ be a semismall proper morphism with $X$ smooth and
  $Y$ irreducible. Then $Rf_*\mathbb{C}_X[\dim X]$ is a semisimple
  perverse sheaf and
  \[
    Rf_*\mathbb{C}_X[\dim X]
    \cong
    \bigoplus_{\alpha}
    \bigoplus_{L\in\operatorname{IrrLoc}(Y_\alpha)}
    IC_{\overline{Y_\alpha}}(L)\otimes V_{\alpha,L}.
  \]
  Only relevant strata occur. For $y\in Y_\alpha$, the local system
  whose isotypic multiplicity spaces are the $V_{\alpha,L}$ has fiber
  $H^{\mathrm{top}}(f^{-1}(y))$.
\end{theorem}

Now consider the Springer resolution
\[
  \pi:\widetilde{\mathcal N}\longrightarrow\mathcal N.
\]
It is semismall by the dimension calculation in
\Cref{sec:steinberg_variety}. Its strata are the nilpotent orbits, and
the monodromy on the top cohomology of a Springer fiber factors through
the component group of the centralizer. The semismall decomposition and
the Springer correspondence therefore give the following statement.

\begin{theorem}[thm:bbd_springer]{\cite{Springer1978,BM83}}
  Let $x\in\mathcal O$ and let
  \(A(x)=\pi_0(Z_G(x))\). Then
  \[
    R\pi_*\mathbb{C}_{\widetilde{\mathcal N}}
      [\dim\widetilde{\mathcal N}]
    \cong
    \bigoplus_{(\mathcal O,\chi)}
    IC_{\overline{\mathcal O}}(L_\chi)
      \otimes\rho_{\mathcal O,\chi},
  \]
  where the sum is over the pairs $(\mathcal O,\chi)$ occurring in the
  Springer correspondence, $L_\chi$ is the local system associated
  with $\chi\in\operatorname{Irr}(A(x))$, and
  $\rho_{\mathcal O,\chi}$ is the corresponding irreducible
  $W$-representation.
\end{theorem}

\subsection{Type \mathintitle{A}}

Take $G=\GL_n$. Nilpotent orbits are indexed by partitions
$\lambda\vdash n$. By \Cref{lem:centr_conn}, every centralizer
$Z_G(x)$ is connected, so $A(x)$ is trivial. Consequently, only the
constant local system occurs on each orbit. Applying the
$T$-equivariant version of the decomposition theorem and taking
hypercohomology gives
\begin{corollary}[crl:type_A_decomposition]{}
  For the type $A$ Springer resolution,
  \[
    \HH_T^\bullet(\widetilde{\mathcal N})
    \cong
    \bigoplus_{\lambda\vdash n}
    \IH_T^\bullet(\overline{\mathcal O}_\lambda)
      \otimes\Htop(\BB_\lambda).
  \]
\end{corollary}

Finally, Lusztig identifies the grading on the intersection-cohomology
factor with Kostka--Foulkes polynomials.
\begin{theorem}[thm:IH_kostka]{\cite{lusztig1981green}}
  For $x\in\mathcal O_\mu$, the Poincar\'e polynomial of the stalk of
  $IC_{\overline{\mathcal O_\lambda}}$ at $x$ is
  \[
    \sum_i q^i
      \dim\mathcal H^i(IC_{\overline{\mathcal O_\lambda}})_x
    =q^{n(\mu)-n(\lambda)}K_{\lambda,\mu}(q^{-1}),
  \]
  where
  \[
    n(\lambda)=\sum_i(i-1)\lambda_i=\dim\BB_\lambda.
  \]
  In particular,
  \[
    Q(\mathcal O_\lambda,q)
    :=\sum_iq^i\dim\IH^i(\overline{\mathcal O_\lambda})
    =q^{\binom n2-n(\lambda)}K_{\lambda,(1^n)}(q^{-1}).
  \]
\end{theorem}
